% Options for packages loaded elsewhere
\PassOptionsToPackage{unicode}{hyperref}
\PassOptionsToPackage{hyphens}{url}
%
\documentclass[
  11pt,
  a4paper,
]{article}
\usepackage{amsmath,amssymb}
\usepackage{iftex}
\ifPDFTeX
  \usepackage[T1]{fontenc}
  \usepackage[utf8]{inputenc}
  \usepackage{textcomp} % provide euro and other symbols
\else % if luatex or xetex
  \usepackage{unicode-math} % this also loads fontspec
  \defaultfontfeatures{Scale=MatchLowercase}
  \defaultfontfeatures[\rmfamily]{Ligatures=TeX,Scale=1}
\fi
\usepackage{lmodern}
\ifPDFTeX\else
  % xetex/luatex font selection
  \setmainfont[]{DejaVu Serif}
  \setsansfont[]{DejaVu Sans}
  \setmonofont[]{DejaVu Sans Mono}
\fi
% Use upquote if available, for straight quotes in verbatim environments
\IfFileExists{upquote.sty}{\usepackage{upquote}}{}
\IfFileExists{microtype.sty}{% use microtype if available
  \usepackage[]{microtype}
  \UseMicrotypeSet[protrusion]{basicmath} % disable protrusion for tt fonts
}{}
\makeatletter
\@ifundefined{KOMAClassName}{% if non-KOMA class
  \IfFileExists{parskip.sty}{%
    \usepackage{parskip}
  }{% else
    \setlength{\parindent}{0pt}
    \setlength{\parskip}{6pt plus 2pt minus 1pt}}
}{% if KOMA class
  \KOMAoptions{parskip=half}}
\makeatother
\usepackage{xcolor}
\usepackage[margin=22mm]{geometry}
\usepackage{color}
\usepackage{fancyvrb}
\newcommand{\VerbBar}{|}
\newcommand{\VERB}{\Verb[commandchars=\\\{\}]}
\DefineVerbatimEnvironment{Highlighting}{Verbatim}{commandchars=\\\{\}}
% Add ',fontsize=\small' for more characters per line
\newenvironment{Shaded}{}{}
\newcommand{\AlertTok}[1]{\textcolor[rgb]{1.00,0.00,0.00}{\textbf{#1}}}
\newcommand{\AnnotationTok}[1]{\textcolor[rgb]{0.38,0.63,0.69}{\textbf{\textit{#1}}}}
\newcommand{\AttributeTok}[1]{\textcolor[rgb]{0.49,0.56,0.16}{#1}}
\newcommand{\BaseNTok}[1]{\textcolor[rgb]{0.25,0.63,0.44}{#1}}
\newcommand{\BuiltInTok}[1]{\textcolor[rgb]{0.00,0.50,0.00}{#1}}
\newcommand{\CharTok}[1]{\textcolor[rgb]{0.25,0.44,0.63}{#1}}
\newcommand{\CommentTok}[1]{\textcolor[rgb]{0.38,0.63,0.69}{\textit{#1}}}
\newcommand{\CommentVarTok}[1]{\textcolor[rgb]{0.38,0.63,0.69}{\textbf{\textit{#1}}}}
\newcommand{\ConstantTok}[1]{\textcolor[rgb]{0.53,0.00,0.00}{#1}}
\newcommand{\ControlFlowTok}[1]{\textcolor[rgb]{0.00,0.44,0.13}{\textbf{#1}}}
\newcommand{\DataTypeTok}[1]{\textcolor[rgb]{0.56,0.13,0.00}{#1}}
\newcommand{\DecValTok}[1]{\textcolor[rgb]{0.25,0.63,0.44}{#1}}
\newcommand{\DocumentationTok}[1]{\textcolor[rgb]{0.73,0.13,0.13}{\textit{#1}}}
\newcommand{\ErrorTok}[1]{\textcolor[rgb]{1.00,0.00,0.00}{\textbf{#1}}}
\newcommand{\ExtensionTok}[1]{#1}
\newcommand{\FloatTok}[1]{\textcolor[rgb]{0.25,0.63,0.44}{#1}}
\newcommand{\FunctionTok}[1]{\textcolor[rgb]{0.02,0.16,0.49}{#1}}
\newcommand{\ImportTok}[1]{\textcolor[rgb]{0.00,0.50,0.00}{\textbf{#1}}}
\newcommand{\InformationTok}[1]{\textcolor[rgb]{0.38,0.63,0.69}{\textbf{\textit{#1}}}}
\newcommand{\KeywordTok}[1]{\textcolor[rgb]{0.00,0.44,0.13}{\textbf{#1}}}
\newcommand{\NormalTok}[1]{#1}
\newcommand{\OperatorTok}[1]{\textcolor[rgb]{0.40,0.40,0.40}{#1}}
\newcommand{\OtherTok}[1]{\textcolor[rgb]{0.00,0.44,0.13}{#1}}
\newcommand{\PreprocessorTok}[1]{\textcolor[rgb]{0.74,0.48,0.00}{#1}}
\newcommand{\RegionMarkerTok}[1]{#1}
\newcommand{\SpecialCharTok}[1]{\textcolor[rgb]{0.25,0.44,0.63}{#1}}
\newcommand{\SpecialStringTok}[1]{\textcolor[rgb]{0.73,0.40,0.53}{#1}}
\newcommand{\StringTok}[1]{\textcolor[rgb]{0.25,0.44,0.63}{#1}}
\newcommand{\VariableTok}[1]{\textcolor[rgb]{0.10,0.09,0.49}{#1}}
\newcommand{\VerbatimStringTok}[1]{\textcolor[rgb]{0.25,0.44,0.63}{#1}}
\newcommand{\WarningTok}[1]{\textcolor[rgb]{0.38,0.63,0.69}{\textbf{\textit{#1}}}}
\usepackage{longtable,booktabs,array}
\usepackage{calc} % for calculating minipage widths
% Correct order of tables after \paragraph or \subparagraph
\usepackage{etoolbox}
\makeatletter
\patchcmd\longtable{\par}{\if@noskipsec\mbox{}\fi\par}{}{}
\makeatother
% Allow footnotes in longtable head/foot
\IfFileExists{footnotehyper.sty}{\usepackage{footnotehyper}}{\usepackage{footnote}}
\makesavenoteenv{longtable}
\setlength{\emergencystretch}{3em} % prevent overfull lines
\providecommand{\tightlist}{%
  \setlength{\itemsep}{0pt}\setlength{\parskip}{0pt}}
\setcounter{secnumdepth}{-\maxdimen} % remove section numbering
\ifLuaTeX
\usepackage[bidi=basic]{babel}
\else
\usepackage[bidi=default]{babel}
\fi
\babelprovide[main,import]{russian}
\ifPDFTeX
\else
\babelfont{rm}[]{DejaVu Serif}
\fi
% get rid of language-specific shorthands (see #6817):
\let\LanguageShortHands\languageshorthands
\def\languageshorthands#1{}
\usepackage{ragged2e}
\usepackage{needspace}
\usepackage{etoolbox}
\usepackage{xurl}
\usepackage{fvextra}
\usepackage{embedfile}
\DefineVerbatimEnvironment{Highlighting}{Verbatim}{breaklines,commandchars=\\\{\},fontsize=\small}
\DefineVerbatimEnvironment{verbatim}{Verbatim}{breaklines,fontsize=\small}
\hfuzz=0.3pt
\AtBeginDocument{%
  \RaggedRight
  \setlength{\emergencystretch}{3em}%
  \hypersetup{pdftitle={LRX: все порядки восьми меток при девяти нулевых блоках}}%
  \embedfile[desc={Nine-gap certificate and standalone verifiers},mimetype={application/zip},ucfilespec={\detokenize{lrx_nine_gap_m8_verification.zip}}]{publication/lrx_nine_gap_m8_verification.zip}%
}
\BeforeBeginEnvironment{longtable}{\Needspace{8\baselineskip}}
\AtBeginEnvironment{longtable}{\small}
\BeforeBeginEnvironment{itemize}{\Needspace{4\baselineskip}}
\pretocmd{\subsection}{\Needspace{7\baselineskip}}{}{}
\pretocmd{\subsubsection}{\Needspace{5\baselineskip}}{}{}
\ifLuaTeX
  \usepackage{selnolig}  % disable illegal ligatures
\fi
\IfFileExists{bookmark.sty}{\usepackage{bookmark}}{\usepackage{hyperref}}
\IfFileExists{xurl.sty}{\usepackage{xurl}}{} % add URL line breaks if available
\urlstyle{same}
\hypersetup{
  pdflang={ru-RU},
  hidelinks,
  pdfcreator={LaTeX via pandoc}}

\author{}
\date{}

\begin{document}

\hypertarget{ux43fux43eux43aux440ux44bux442ux438ux435-lrx-ux432ux441ux435-ux43fux43eux440ux44fux434ux43aux438-ux432ux43eux441ux44cux43cux438-ux43cux435ux442ux43eux43a-ux43fux440ux438-ux434ux435ux432ux44fux442ux438-ux43dux443ux43bux435ux432ux44bux445-ux431ux43bux43eux43aux430ux445}{%
\section{Покрытие LRX: все порядки восьми меток при девяти нулевых
блоках}\label{ux43fux43eux43aux440ux44bux442ux438ux435-lrx-ux432ux441ux435-ux43fux43eux440ux44fux434ux43aux438-ux432ux43eux441ux44cux43cux438-ux43cux435ux442ux43eux43a-ux43fux440ux438-ux434ux435ux432ux44fux442ux438-ux43dux443ux43bux435ux432ux44bux445-ux431ux43bux43eux43aux430ux445}}

24 сентября 2026 года. Доказательство частного случая и перенос
сертификатов удалением нулевых блоков.

\hypertarget{ux442ux435ux43eux440ux435ux43cux430-ux438-ux43eux431ux43bux430ux441ux442ux44c-ux440ux435ux437ux443ux43bux44cux442ux430ux442ux430}{%
\subsection{1. Теорема и область
результата}\label{ux442ux435ux43eux440ux435ux43cux430-ux438-ux43eux431ux43bux430ux441ux442ux44c-ux440ux435ux437ux443ux43bux44cux442ux430ux442ux430}}

Операции \(L,R\) циклически сдвигают вектор влево и вправо на одну
позицию, \(X\) меняет первые два элемента. Каждая операция стоит один
шаг. Через \(d(v)\) обозначим расстояние до корня \((1,\ldots,8,0^r)\).

\textbf{Теорема 1.} Для любой перестановки \((a_1,\ldots,a_8)\) чисел
\(1,\ldots,8\) и любых положительных целых \(\ell_0,\ldots,\ell_8\)
вектор

\[v=(0^{\ell_0},a_1,0^{\ell_1},a_2,\ldots,a_8,0^{\ell_8})\]

удовлетворяет

\[\boxed{d(v)\le6n-18
=\binom n2-\frac{(r-1)(r+4)}2,}
\qquad r=\sum_{j=0}^{8}\ell_j,\quad n=8+r.\tag{1}\]

Иными словами, нулевой блок имеется в каждом из девяти линейных
промежутков, включая оба конца отображаемого вектора. Порядок меток
произволен. Длины блоков не ограничены сверху.

Теорема компьютерно-ассистированная: общие леммы ниже сводят её к
конечному набору точных рациональных неравенств. Для всех \(8!=40320\)
порядков эти неравенства независимо проверены. Сортирующие слова даны
явно в автономном пакете данных.

Удаление части нулевых блоков дополнительно даёт критерий, доказавший ту
же границу для 12680558 базовых семейств при \(m=8\). Здесь базовое
семейство задаётся порядком восьми меток и непустым набором нулевых
промежутков; длины всех оставшихся блоков произвольны и положительны. С
учётом ранее доказанных классов покрыто 16107991 из 20603520 таких
семейств. Остались 4495529 семейств с 4--8 блоками.

\textbf{Общая гипотеза для всех \(m=n-r\ge8\) не доказана.} Также пока
не доказан полный случай \(m=8\). Результат относится к покрытию из
канонического корня; диаметр между произвольными вершинами здесь не
рассматривается.

\hypertarget{ux441ux440ux430ux432ux43dux435ux43dux438ux44f-ux432ux43cux435ux441ux442ux43e-ux431ux435ux437ux443ux441ux43bux43eux432ux43dux44bux445-ux43eux431ux43cux435ux43dux43eux432}{%
\subsection{2. Сравнения вместо безусловных
обменов}\label{ux441ux440ux430ux432ux43dux435ux43dux438ux44f-ux432ux43cux435ux441ux442ux43e-ux431ux435ux437ux443ux441ux43bux43eux432ux43dux44bux445-ux43eux431ux43cux435ux43dux43eux432}}

Закрепим физические позиции окружности и подвижный курсор. \(L\)
увеличивает координату курсора, \(R\) уменьшает её, а \(X\) меняет
элементы на его позиции и следующей позиции. Видимый вектор начинается у
курсора.

Пусть слово \(W\) сортирует опорный вектор \(q\) длины \(N\) и
заканчивает с курсором \(f\). Выберем разрез перед позицией \(c\), через
который ни один обмен \(X\) не проходит. Конечная физическая цель метки
\(x\) имеет координату \((f+x-1)\bmod N\).

Развернём окружность в окно, начинающееся в \(c\). Присвоим элементам
ранги их целевых позиций \(0,\ldots,N-1\) в этом окне. Нули сопоставим
нулевым целям по возрастанию. Если \(W\) не обменивает два нуля, их
линейный порядок сохраняется. Таким образом, \(W\) сортирует полученную
перестановку рангов \(Q\).

Для другого входа \(p\) с теми же нулевыми позициями обозначим его
перестановку целевых рангов через \(P\). Положим

\[H_P(j,s)=\#\{i<j:P_i<s\}.\]

Условие префиксного доминирования записываем в виде

\[P\preceq Q
\quad\Longleftrightarrow\quad
H_P(j,s)\ge H_Q(j,s)\quad\text{для всех }j,s.\tag{2}\]

\textbf{Лемма 2.} При (2) то же слово позиций обмена сортирует \(P\),
если каждый обмен выполнять лишь тогда, когда пара рангов убывает. Все
исходные вращения можно сохранить. Выполненных обменов будет ровно
\(\operatorname{inv}(P)\).

\textbf{Доказательство.} Для фиксированного порога \(s\) сравнение на
позициях \(j-1,j\) меняет только префиксное число длины \(j\); его новое
значение равно

\[\min\bigl(H_P(j-1,s)+1,H_P(j+1,s)\bigr).\tag{3}\]

Это максимальное значение, достижимое перестановкой данной пары. Функция
(3) не убывает по обоим аргументам. Поэтому после сравнения на \(P\) и
безусловного обмена на \(Q\) условие (2) сохраняется. В конце \(Q\)
отсортирована и имеет максимально возможные префиксные числа
\(\min(j,s)\). Значит, \(P\) также отсортирована. Каждый выполненный
обмен уменьшает число инверсий на один. \(\square\)

Поскольку вращения сохранены, конечный курсор по-прежнему равен \(f\).
Видимый конечный вектор точно равен каноническому корню. В этой
аргументации нет бесплатной смены фазы.

Условие (2) проверяется конечным способом: для каждой длины префикса
отсортировать его ранги и потребовать, чтобы каждый ранг \(P\) был не
больше соответствующего ранга \(Q\). Опорный порядок меток может быть
произвольным; обратный порядок не предполагается.

\hypertarget{ux440ux430ux441ux442ux44fux436ux435ux43dux438ux435-ux438-ux442ux43eux447ux43dux430ux44f-ux446ux435ux43dux430-ux441ux440ux430ux432ux43dux438ux432ux430ux44eux449ux435ux433ux43e-ux441ux43bux43eux432ux430}{%
\subsection{3. Растяжение и точная цена сравнивающего
слова}\label{ux440ux430ux441ux442ux44fux436ux435ux43dux438ux435-ux438-ux442ux43eux447ux43dux430ux44f-ux446ux435ux43dux430-ux441ux440ux430ux432ux43dux438ux432ux430ux44eux449ux435ux433ux43e-ux441ux43bux43eux432ux430}}

Сначала все нулевые атомы опорного вектора одиночные. Растянем \(j\)-й
атом до длины \(\ell_j=1+z_j\), где \(z_j\ge0\). Сортировка переносится
следующими буквальными макросами:

\begin{longtable}[]{@{}ll@{}}
\toprule\noalign{}
Атомарное действие & Действие после растяжения \\
\midrule\noalign{}
\endhead
\bottomrule\noalign{}
\endlastfoot
\(L\) через атом длины \(u\) & \(L^u\) \\
\(R\) через атом длины \(u\) & \(R^u\) \\
\(X\) на \((0^u,x)\) & \(L^{u-1}X(RX)^{u-1}\) \\
\(X\) на \((x,0^u)\) & \(X(LX)^{u-1}R^{u-1}\) \\
\(X\) двух одиночных меток & \(X\) \\
\end{longtable}

Каждая строка выполняет указанное действие над целыми атомами. После
последнего макроса видимый вектор равен точному корню. Если исходный
обмен не пересекал выбранный разрез, соответствующий макрос также его не
пересекает: он действует внутри двух соседних атомов. Разрез остаётся
границей атомов.

\hypertarget{ux446ux435ux43dux430-ux43eux43fux43eux440ux43dux43eux433ux43e-ux43fux43eux434ux44aux451ux43cux430}{%
\subsubsection{Цена опорного
подъёма}\label{ux446ux435ux43dux430-ux43eux43fux43eux440ux43dux43eux433ux43e-ux43fux43eux434ux44aux451ux43cux430}}

Пусть \(q_X\) --- число букв \(X\) исходного слова, \(s_j\) --- число
обменов \(j\)-го нуля с метками. После выделения внутренних служб в
макросах между ними остаются вращательные промежутки со знаковыми
длинами \(c_g+\sum_j a_{gj}z_j\). Отсюда

\[\begin{aligned}
|W(\ell)|&=q_X+2\sum_js_jz_j+
\sum_g\left|c_g+\sum_ja_{gj}z_j\right|\\
&\le B+\sum_j\beta_jz_j,\\
B&=q_X+\sum_g|c_g|,\qquad
\beta_j=2s_j+\sum_g|a_{gj}|.
\end{aligned}\tag{4}\]

Все коэффициенты находятся проходом по слову с помеченными нулями. В
используемом пакете \(B=|W|\), а постоянная часть и коэффициенты каждого
промежутка имеют один знак. Поэтому верхняя цена (4) точна при всех
\(z_j\ge0\). Эти условия проверяются заново для каждого слова.

\hypertarget{ux44dux43aux43eux43dux43eux43cux438ux44f-ux43eux431ux43cux435ux43dux43eux432}{%
\subsubsection{Экономия
обменов}\label{ux44dux43aux43eux43dux43eux43cux438ux44f-ux43eux431ux43cux435ux43dux43eux432}}

Предположим, что \(q_X=\operatorname{inv}(Q)\). Тогда все обмены
опорного слова уменьшают число инверсий. Через \(\sigma_j(P)\) обозначим
число инверсий \(j\)-го нуля с метками в \(P\). Положим

\[\delta_0=\operatorname{inv}(Q)-\operatorname{inv}(P),\qquad
\delta_j=\sigma_j(Q)-\sigma_j(P).\tag{5}\]

При (2) все эти величины неотрицательны. Для \(\delta_0\) это следует из
леммы 2: сравнение использует \(\operatorname{inv}(P)\) обменов из
имеющихся \(\operatorname{inv}(Q)\) позиций обмена.

Для нуля пусть перед ним стоят \(h\) меток и ровно \(s\) меток имеют
целевой ранг меньше его ранга. Пусть \(p,q\) обозначают
последовательности сжатых рангов меток в \(P,Q\) соответственно. Тогда

\[\sigma_j(P)=h+s-2H_p(h,s).\]

Следовательно, \(\delta_j=2(H_p(h,s)-H_q(h,s))\ge0\). В частности,
каждое \(\delta_j\) чётно.

Нулевые ранги и их исходные позиции одинаковы для \(P,Q\). В условии (2)
они сокращаются, и остаются только префиксные условия для меток. При
растяжении порядок меток и их целевых рангов сохраняется; добавленные
нули снова вносят одинаковые слагаемые. Поэтому (2) выполняется при всех
положительных длинах.

Каждый дополнительный экземпляр нуля имеет те же отношения порядка с
метками, что исходный экземпляр. Отсюда

\[\operatorname{inv}(P_\ell)=\operatorname{inv}(P)
+\sum_j\sigma_j(P)z_j.\tag{6}\]

Для опорного подъёма число обменов также равно
\(\operatorname{inv}(Q_\ell)\): каждый первоначальный обмен нуля с
меткой заменяется на \(\ell_j\) обменов, а порядок нулей сохраняется.

\textbf{Лемма 3.} Сначала растянем опорное слово, затем выполним
сравнения по лемме 2, сохраняя все опорные вращения. Длина полученного
слова равна

\[\boxed{C_{P,W}(\ell)=B-\delta_0+
\sum_j(\beta_j-\delta_j)(\ell_j-1).}\tag{7}\]

\textbf{Доказательство.} Вращения одинаковы в двух исполнениях. Разность
чисел обменов равна
\(\operatorname{inv}(Q_\ell)-\operatorname{inv}(P_\ell) =\delta_0+\sum_j\delta_jz_j\)
по (6). Вычитание из (4) даёт (7). \(\square\)

Равенство относится к построенному слову с сохранёнными вращениями, а не
к кратчайшему расстоянию. Дополнительные сокращения могут сделать его
короче.

\hypertarget{ux440ux430ux446ux438ux43eux43dux430ux43bux44cux43dux44bux439-ux441ux435ux440ux442ux438ux444ux438ux43aux430ux442-ux432ux441ux435ux445-ux434ux43bux438ux43d}{%
\subsection{4. Рациональный сертификат всех
длин}\label{ux440ux430ux446ux438ux43eux43dux430ux43bux44cux43dux44bux439-ux441ux435ux440ux442ux438ux444ux438ux43aux430ux442-ux432ux441ux435ux445-ux434ux43bux438ux43d}}

Зафиксируем порядок \(a=(a_1,\ldots,a_8)\). Его единичная база имеет
длину 17 и девять нулей. Каждая строка сертификата задаёт опорную базу с
теми же нулевыми позициями, слово \(W\), допустимый разрез \(c\) и
положительный рациональный вес \(\theta_W\).

Для каждой строки проверяются: сортировка опорной базы; отсутствие
обменов через разрез; условие (2) для данного \(a\);
\(q_X=\operatorname{inv}(Q)\); формула (4) и величины (5). Обозначим
базовую цену и коэффициенты (7) через \(b_W,\gamma_{Wj}\).

\textbf{Лемма 4.} Если

\[\sum_W\theta_W=1,\qquad
\sum_W\theta_W b_W<85,\qquad
\sum_W\theta_W\gamma_{Wj}\le6\quad(0\le j\le8),\tag{8}\]

то (1) верно при всех \(\ell_j\ge1\).

\textbf{Доказательство.} Средняя цена не превосходит

\[\sum_W\theta_Wb_W+6\sum_j(\ell_j-1)
<31+6\sum_j\ell_j=(6n-18)+1.\]

Одно из слов имеет длину не выше средней. Его длина целая, поэтому она
не больше \(6n-18\). \(\square\)

Таким образом, бесконечное множество длин проверяется конечным набором
рациональных неравенств. Итоговый проверяющий читает веса точными
дробями и заново проверяет (8); оптимизатор в нём не используется.

\hypertarget{ux43fux43eux43bux43dux43eux442ux430-ux43aux43eux43dux435ux447ux43dux43eux433ux43e-ux441ux43fux438ux441ux43aux430}{%
\subsubsection{Полнота конечного
списка}\label{ux43fux43eux43bux43dux43eux442ux430-ux43aux43eux43dux435ux447ux43dux43eux433ux43e-ux441ux43fux438ux441ux43aux430}}

\href{literature/multiset_nine_gap_complete_certificates_20260924.json}{Пакет
сертификатов} содержит каталог опорных слов и ровно 40320 записей.
Запись задаёт порядок меток и список троек «номер опорного слова,
разрез, вес». Проверяющий самостоятельно перечисляет все перестановки
восьми меток и требует точного совпадения множеств, без пропусков и
повторов.

Для каждой записи проверены (8) и все предпосылки лемм 2--3.
Следовательно, лемма 4 применима ко всем порядкам меток. Это завершает
доказательство теоремы 1. \(\square\)

\hypertarget{ux43eux431ux44aux451ux43c-ux43dux435ux437ux430ux432ux438ux441ux438ux43cux43eux439-ux43fux440ux43eux432ux435ux440ux43aux438}{%
\subsection{5. Объём независимой
проверки}\label{ux43eux431ux44aux451ux43c-ux43dux435ux437ux430ux432ux438ux441ux438ux43cux43eux439-ux43fux440ux43eux432ux435ux440ux43aux438}}

\href{scripts/audit_multiset_nine_gap_cover.py}{Проверяющий} не
импортирует оптимизатор, генератор слов или механизм вычисления областей
переноса. Он независимо восстанавливает физическую цель, ранги и
префиксные условия.

\begin{longtable}[]{@{}lr@{}}
\toprule\noalign{}
Проверенный объект & Количество \\
\midrule\noalign{}
\endhead
\bottomrule\noalign{}
\endlastfoot
Порядки меток, без пропусков & 40320 \\
Опорные слова & 4670 \\
Строки рациональных смесей & 152937 \\
Буквальные проверки коэффициентов растяжения & 42030 \\
Сравнивающие слова на единичных базах & 40320 \\
Буквы в единичных исполнениях & 2859789 \\
Дополнительные расширенные слова & 1260 \\
Буквы в расширенных исполнениях & 261839 \\
\end{longtable}

Проверяющий отверг три намеренно испорченных случая: неверный вес,
разрез с обменом через него и неполный список порядков.
\href{literature/multiset_nine_gap_complete_audit_20260924.json}{Основной
аудит} собрал автономный пакет;
\href{literature/multiset_nine_gap_standalone_audit_20260924.json}{повторный
аудит} прочитал только этот пакет и получил те же результаты.

Конечные буквальные исполнения проверяют реализацию. Обоснование
неограниченных длин --- леммы 2--4, а не экстраполяция этих исполнений.
Доказательство не использует спорную формулировку \(2\Delta+K^2\le b^2\)
или SMT-приложения из \texttt{input}.

\hypertarget{ux443ux434ux430ux43bux435ux43dux438ux435-ux431ux43bux43eux43aux43eux432-ux43eux431ux449ux430ux44f-ux43bux435ux43cux43cux430-ux438-ux434ux43eux43fux43eux43bux43dux438ux442ux435ux43bux44cux43dux44bux435-ux441ux435ux43cux435ux439ux441ux442ux432ux430}{%
\subsection{6. Удаление блоков: общая лемма и дополнительные
семейства}\label{ux443ux434ux430ux43bux435ux43dux438ux435-ux431ux43bux43eux43aux43eux432-ux43eux431ux449ux430ux44f-ux43bux435ux43cux43cux430-ux438-ux434ux43eux43fux43eux43bux43dux438ux442ux435ux43bux44cux43dux44bux435-ux441ux435ux43cux435ux439ux441ux442ux432ux430}}

\hypertarget{ux43aux43eux43cux43cux443ux442ux438ux440ux43eux432ux430ux43dux438ux435-ux443ux434ux430ux43bux435ux43dux438ux44f-ux441-ux440ux430ux441ux442ux44fux436ux435ux43dux438ux435ux43c}{%
\subsubsection{Коммутирование удаления с
растяжением}\label{ux43aux43eux43cux43cux443ux442ux438ux440ux43eux432ux430ux43dux438ux435-ux443ux434ux430ux43bux435ux43dux438ux44f-ux441-ux440ux430ux441ux442ux44fux436ux435ux43dux438ux435ux43c}}

Пометим исходные нули. Выберем множество удаляемых нулевых атомов \(D\);
хотя бы один нуль оставим. Проекция отдельного шага устроена так:

\begin{itemize}
\tightlist
\item
  \(L\) сохраняется, если проходимый первый атом оставлен;
\item
  \(R\) сохраняется, если проходимый последний атом оставлен;
\item
  \(X\) сохраняется, если оба обмениваемых атома оставлены;
\item
  в остальных случаях шаг исчезает.
\end{itemize}

Во всех четырёх случаях непосредственная проверка показывает: сначала
выполнить исходный шаг и удалить помеченные атомы --- то же действие над
видимым вектором, что сначала удалить их и выполнить проекцию шага.
Индукция по слову даёт сортирующее слово \(W^D\) на укороченной базе.
Его начальная и конечная ориентации сохраняются по этой пошаговой
проверке.

Удаляемые атомы оставим единичными, а остальные растянем произвольно.
Для макросов §3 действует то же тождество. Макрос вращения через
удаляемый атом исчезает. Макрос обмена с таким нулём при его длине один
состоит из единственного \(X\) и также исчезает. Макросы на оставленных
атомах сохраняются целиком. Следовательно,

\[\operatorname{erase}_D\bigl(\operatorname{lift}(W)\bigr)
=\operatorname{lift}\bigl(\operatorname{erase}_D(W)\bigr),\tag{9}\]

с точностью до сокращения соседних \(LR,RL\). Эти сокращения тоже
коммутируют с удалением: обратная пара проходит один и тот же атом,
поэтому обе её буквы либо сохраняются, либо исчезают.

\textbf{Лемма 5.} Для коэффициентов стандартного подъёма (4) удаление
других нулевых атомов даёт

\[\beta_j^D\le\beta_j\qquad(j\notin D).\tag{10}\]

\textbf{Доказательство.} Проекция только удаляет буквы и может сократить
обратные вращения. По (9) длина поднятого проекционного слова не больше
длины исходного подъёма при любых оставленных длинах. Увеличим только
\(\ell_j\), остальные положим равными единице. Старшая линейная часть
точной формулы с модулями в (4) имеет коэффициент \(\beta_j\); для
проекции --- \(\beta_j^D\). Неравенство длин при сколь угодно больших
целых \(\ell_j\) даёт (10). \(\square\)

\hypertarget{ux441ux43eux445ux440ux430ux43dux435ux43dux438ux435-ux443ux441ux43bux43eux432ux438ux44f-ux441ux440ux430ux432ux43dux435ux43dux438ux44f}{%
\subsubsection{Сохранение условия
сравнения}\label{ux441ux43eux445ux440ux430ux43dux435ux43dux438ux435-ux443ux441ux43bux43eux432ux438ux44f-ux441ux440ux430ux432ux43dux435ux43dux438ux44f}}

В перестановках \(P,Q\) удаляемый нуль имеет одинаковые исходную позицию
и целевой ранг. При удалении такого общего элемента из префиксных чисел
вычитается одинаковый вклад. Поэтому (2) сохраняется. Разрез переходит в
границу перед первым оставленным атомом после него. Ни один оставленный
обмен эту границу не пересекает.

Удаление других нулей не меняет инверсий оставленного нуля с метками.
Следовательно, \(\delta_j\) для \(j\notin D\) сохраняется, а

\[\delta_0^D=\delta_0-\sum_{j\in D}\delta_j.\tag{11}\]

Каждый обмен опорного слова пересекал инвертированную пару ровно один
раз. После удаления нулей это остаётся верным для оставшихся пар.
Значит, спроецированное опорное слово также имеет минимальное число
обменов в своём окне, и к нему применима лемма 3.

Пусть \(B_W^D\) --- буквально посчитанная длина спроецированного
опорного слова после сокращения \(LR,RL\). Его новая базовая цена для
данного входа равна

\[b_W^D=B_W^D-\delta_0+\sum_{j\in D}\delta_j.\tag{12}\]

Новые коэффициенты сравнивающего подъёма не больше старых:
\(\beta_j^D-\delta_j\le\beta_j-\delta_j\) по (10). Поэтому прежние веса
сохраняют ограничения на все оставшиеся коэффициенты. Если осталось
\(k\) блоков, для нового сертификата достаточно проверить единственное
дополнительное неравенство

\[\sum_W\theta_W b_W^D<31+6k.\tag{13}\]

Именно (13), а не само существование проекции, оплачивает уменьшение
бюджета на шесть шагов за каждый удалённый единичный нуль.

\hypertarget{ux43fux43eux43bux43dux430ux44f-ux43fux440ux43eux432ux435ux440ux43aux430-ux432ux441ux435ux445-ux43fux43eux434ux43cux43dux43eux436ux435ux441ux442ux432}{%
\subsubsection{Полная проверка всех
подмножеств}\label{ux43fux43eux43bux43dux430ux44f-ux43fux440ux43eux432ux435ux440ux43aux430-ux432ux441ux435ux445-ux43fux43eux434ux43cux43dux43eux436ux435ux441ux442ux432}}

Для каждого из 4670 опорных слов рассмотрены все 511 непустых множеств
оставленных нулей. Получены 2386370 проекций.
\href{scripts/verify_multiset_projection_table.cpp}{Второй проверяющий
на C++} не использует код построения таблицы: он буквально вращает
списки, исполняет проекцию и сверяет её конечный вектор и длину. Все
цены совпали.

\href{literature/multiset_nine_gap_projection_audit_20260924.json}{Аудит
проекционных неравенств} заново пересчитал все
\(40320\cdot511=20603520\) экземпляров (13). Веса приведены к общему
знаменателю; целочисленные вычисления имеют проверенную границу,
исключающую переполнение. Для 5110 экземпляров выполнена дополнительная
сверка скалярными точными дробями. Намеренно неверная цена проекции
отвергнута.

\begin{longtable}[]{@{}rrr@{}}
\toprule\noalign{}
Оставшихся блоков & Всего базовых семейств & Покрыто проекцией \\
\midrule\noalign{}
\endhead
\bottomrule\noalign{}
\endlastfoot
1 & 362880 & 81255 \\
2 & 1451520 & 448700 \\
3 & 3386880 & 1415077 \\
4 & 5080320 & 2783585 \\
5 & 5080320 & 3484403 \\
6 & 3386880 & 2752058 \\
7 & 1451520 & 1322558 \\
8 & 362880 & 352602 \\
9 & 40320 & 40320 \\
\textbf{Всего} & \textbf{20603520} & \textbf{12680558} \\
\end{longtable}

Объединение с \href{LRX_MULTISET_THREE_BLOCK_M8_THEOREM_RU.md}{теоремой
о не более трёх блоках} и
\href{LRX_MULTISET_COMPARISON_TRANSFER_RU.md}{прежним переносом обратных
порядков} содержит 16107991 семейство. Эта цифра использует указанные
ранее доказанные результаты как отдельные основания.

\hypertarget{ux43fux43eux447ux435ux43cux443-ux43eux434ux43dux43eux433ux43e-ux440ux430ux441ux442ux44fux433ux438ux432ux430ux435ux43cux43eux433ux43e-ux441ux43bux43eux432ux430-ux43dux435ux434ux43eux441ux442ux430ux442ux43eux447ux43dux43e}{%
\subsection{7. Почему одного растягиваемого слова
недостаточно}\label{ux43fux43eux447ux435ux43cux443-ux43eux434ux43dux43eux433ux43e-ux440ux430ux441ux442ux44fux433ux438ux432ux430ux435ux43cux43eux433ux43e-ux441ux43bux43eux432ux430-ux43dux435ux434ux43eux441ux442ux430ux442ux43eux447ux43dux43e}}

Следующая лемма справедлива для любого числа одиночных меток \(m\ge2\).
Она относится к методу фиксированного макроподъёма и не опровергает
гипотезу о расстояниях.

Пусть между любыми двумя соседними по окружности метками имеется хотя бы
один нулевой атом. Рассмотрим слово без обменов двух нулей со
стандартным макроподъёмом §3 и закреплёнными пометками атомов. Различим
нули и занумеруем их в круговом порядке. Этот порядок сохраняется.
Свободная смена пометок внутри слившегося нулевого блока в
рассматриваемую конструкцию не входит.

Пусть \(c_j\) --- начальное число меток между нулём \(j\) и следующим
нулём. По условию \(c_j\le1\) и \(\sum_jc_j=m\). В конце все метки
находятся в одном промежутке, скажем \(g\); поэтому \(c'_g=m\),
\(c'_j=0\) при \(j\ne g\).

Обозначим через \(F_j\) знаковое число проходов меток через нуль \(j\),
положительное при движении от предыдущего промежутка к следующему. Через
\(s_j\) обозначим полное число таких обменов. Тогда

\[c'_j-c_j=F_j-F_{j+1},\qquad |F_j|\le s_j.\tag{14}\]

Разрежем круг индексов после промежутка \(g\). На остальных промежутках
потоки последовательно увеличиваются на \(c_j\ge0\). Суммарное
увеличение равно \(m-c_g\). Поэтому

\[\max_jF_j-\min_jF_j=m-c_g\ge m-1.\tag{15}\]

Следовательно, хотя бы один \(s_j\) не меньше
\(\lceil(m-1)/2\rceil=\lfloor m/2\rfloor\). Из формулы (4) получаем

\[\boxed{\max_j\beta_j\ge2\lfloor m/2\rfloor>m-2.}\tag{16}\]

Такой фиксированный подъём не может при всех независимых длинах иметь
цену с коэффициентами не выше \(m-2\): достаточно увеличивать длину
блока, на котором (16) нарушает требуемый коэффициент. При \(m=8\) хотя
бы один коэффициент не меньше восьми, что превышает требуемый предел
шесть.

Препятствие сохраняется для одиночного сравнивающего подъёма (7). При
сравнении \(s_j^{\rm new}\le s_j^{\rm old}\), а его коэффициент не
меньше

\[2s_j^{\rm old}-(s_j^{\rm old}-s_j^{\rm new})
=s_j^{\rm old}+s_j^{\rm new}\ge2s_j^{\rm new}.\]

К новому единичному слову снова применимы (14)--(15). Рациональные смеси
обходят это ограничение: разные слова могут иметь большой коэффициент на
разных блоках; для оценки расстояния выбирается лучшее слово при данных
длинах.

\href{literature/multiset_zero_flux_audit_20260924.json}{Буквальный
аудит потоков} проверил 2166 слов, 10000 схем заполнения промежутков и
15444 проверки точной целочисленной границы. В
\href{lean/MultisetZeroFlux.lean}{Lean-файле} доказаны телескопическая
формула и четыре скалярных следствия. Связь потоков с LRX-словами
доказана выше; полной формализации модели графа в этом файле нет.

\newpage

\hypertarget{ux432ux43eux441ux43fux440ux43eux438ux437ux432ux435ux434ux435ux43dux438ux435-ux438-ux43eux441ux442ux430ux432ux448ux430ux44fux441ux44f-ux437ux430ux434ux430ux447ux430}{%
\subsection{8. Воспроизведение и оставшаяся
задача}\label{ux432ux43eux441ux43fux440ux43eux438ux437ux432ux435ux434ux435ux43dux438ux435-ux438-ux43eux441ux442ux430ux432ux448ux430ux44fux441ux44f-ux437ux430ux434ux430ux447ux430}}

Автономная проверка теоремы 1 из корня проекта:

\begin{Shaded}
\begin{Highlighting}[]
\ExtensionTok{python}\NormalTok{ scripts/audit\_multiset\_nine\_gap\_cover.py }\DataTypeTok{\textbackslash{}}
  \AttributeTok{{-}{-}bundle{-}input}\NormalTok{ literature/multiset\_nine\_gap\_complete\_certificates\_20260924.json }\DataTypeTok{\textbackslash{}}
  \AttributeTok{{-}{-}output}\NormalTok{ /tmp/lrx\_nine\_gap\_recheck.json}
\end{Highlighting}
\end{Shaded}

Для проекционного дополнения:

\begin{Shaded}
\begin{Highlighting}[]
\ExtensionTok{g++} \AttributeTok{{-}O3} \AttributeTok{{-}std}\OperatorTok{=}\NormalTok{c++17 }\AttributeTok{{-}Wall} \AttributeTok{{-}Wextra} \AttributeTok{{-}Werror} \DataTypeTok{\textbackslash{}}
\NormalTok{  scripts/verify\_multiset\_projection\_table.cpp }\DataTypeTok{\textbackslash{}}
  \AttributeTok{{-}o}\NormalTok{ /tmp/lrx\_verify\_projection\_table}
\ExtensionTok{python}\NormalTok{ scripts/audit\_multiset\_nine\_gap\_projections.py }\DataTypeTok{\textbackslash{}}
  \AttributeTok{{-}{-}binary}\NormalTok{ /tmp/lrx\_verify\_projection\_table }\DataTypeTok{\textbackslash{}}
  \AttributeTok{{-}{-}output}\NormalTok{ /tmp/lrx\_nine\_gap\_projection\_recheck.json}
\end{Highlighting}
\end{Shaded}

Проверки следует запускать без \texttt{-O} у Python и с новыми именами
выходных JSON. Первый сценарий не требует SciPy или C++. Во втором NumPy
используется для конечных целочисленных таблиц, а не для приближённого
решения неравенств.

При \(m=8\) остаются семейства с 4--8 блоками, не прошедшие (13) и
прежние критерии. Для них можно менять смеси после проекции или строить
новые опорные слова. Отдельно требуется общий аргумент для произвольного
\(m>8\). Ни полнота девятиблочного случая, ни число покрытых базовых
семейств этот последний шаг не заменяют.

\end{document}
